\chapter{NQR Mail-Screening System Study}
\label{chap:nqr-screening}

A successful spin simulation is only one part of an NQR screening instrument.
The signal must survive spatially varying excitation, polarization transport,
coil loading, receiver dynamics and interference before a calibrated detector
can identify it. The versioned study in
\code{studies/nqr\_mail\_screening} composes these layers with explicit physical
time and retained spin histories. These scripts are reproducible research
entry points, not a new installed package API or a validated screening product.

\chapterpath{Locate the staged reference calculations, understand their
numerical gates, and distinguish a matched synthetic detector test from
demonstrated physical sensitivity.}

\section{Implemented Stages}

\begin{description}
\item[Phase 0: requirements and evidence.] Versioned target/material inputs,
requirements, literature provenance and a frozen test-set policy.
\item[Phase 1: absolute reference.] Three-level pulse propagation and an
absolute signal-chain reference, including coil geometry, thermal populations,
reciprocal reception, noise and convergence checks. Motion-aware pulse histories
and a Zeeman-spectral spacing constraint replace an instantaneous transfer gain.
\item[Phase 2: spatial and material loading.] Voxel-resolved transmit/receive
coupling, circuit parameters, conservative loading bounds and an independent
quasistatic material audit. Validation-only dependencies are isolated in
the phase's requirements file. This is not a general full-wave validation.
\item[Phase 3: complete acquisition records.] Spin density matrices survive
successive SLSE-x, FID-y, SORC-x and repeated-SLSE blocks. Carrier-frame handoffs,
finite RF tails, receiver filtering, ADC clipping/recovery, RFI cancellation and
thermal accounting share one physical clock. Magnet and coil dwells are
stationary, with explicit transfer motion between them.
\item[Phase 4: calibrated decisions.] A three-action adaptive tree, matched-filter
features and correlated likelihood compare fixed and adaptive schedules under
the same complete cycle-time budget. Development, threshold calibration,
held-out evaluation and shifted-interference data use separate records.
\end{description}

\section{Results and Boundaries}

The checked-in Phase 4 report's physical 1 g reference gives ROC AUC 0.509
for the fixed schedule and 0.468 for the adaptive schedule at a 3.354 s budget.
The reported intervals overlap chance, and the paired comparison does not
establish adaptive benefit. Software amplitude multipliers of 10000 and 100000
are explicitly diagnostic positive controls, not extra hardware gain or an
achievable sensitivity claim. Untrained RFI drift raises false alarms far above
the illustrative 5\% matched-domain calibration setting.

Numerical gates therefore establish their declared synthetic domains, not
fentanyl sensitivity, installed-site robustness, blind benign-material
discrimination or optimal hardware. Subsequent work must address physical
signal strength and robustness before site calibration. The study's report
files and phase READMEs preserve these restrictions and provenance.

\section{Reproduce Without Confusing Studies and APIs}

From \code{PythonSpinDynamics}, first read the study README and the chosen
phase's prerequisites. Full spatial and receiver stages can be expensive;
they should not be treated as installation smoke examples.
\begin{lstlisting}[language=bash]
python -m unittest tests.test_nqr_mail_screening_phase4
# After generating the Phase 2 maps and Phase 3 dependencies:
python studies/nqr_mail_screening/phase4/evaluate.py
python studies/nqr_mail_screening/phase4/plot_phase4.py
\end{lstlisting}

The underlying public physics, field, receiver and Bayesian-design APIs are
documented in their respective chapters. Phase scripts, their schemas and
large cached artifacts are study-specific. Numerical test success does not
promote them to a general production workflow.
